% GENERATED FILE. Edit canonical lessons, then run npm run generate:books.
% canonical-source-hash: fnv1a64:4a35a5e72b90dddc
% canonical-lessons: MR-C351-matra, MR-C351-tasec, MR-C351-kimbahuna, MR-C351-tathapi, MR-C351-udaharanartha

\chapter{Linking Ideas}
\label{ch:mr-linking-ideas}
\hlchaptermodality{351}

\begin{chapteropening}
\textbf{By the end of this chapter:} \emph{I can join ideas with however, likewise and for example.}

Complete the last lesson of chapter 351: I can join ideas with however, likewise and for example.
\end{chapteropening}

\section[\texorpdfstring{\mr{मात्र} (mātra)}{मात्र (mātra)}]{\mr{मात्र} (mātra) — however; though}

\label{lesson:MR-C351-matra}



\begin{quote}
Before the new one: say the Marathi for to blame, then the Marathi for to accuse.
\end{quote}

\subsection*{You'll want to know: \mr{मात्र}}
\textbf{\mr{मात्र}} — \emph{mātra} — \textquotedblleft{}however; though\textquotedblright{}.

\begin{quote}
\mr{घर} \mr{लहान} \mr{आहे}, \mr{मात्र} \mr{खूप} \mr{स्वच्छ} \mr{आहे}.
\end{quote}

\begin{quote}
\emph{ghar lahān āhe, mātra khūp svaccha āhe.}
\end{quote}

\begin{quote}
The house is small; it is very clean, though.
\end{quote}

Often after the first word: \mr{मी} \mr{मात्र} \mr{घरीच} \mr{राहिलो}, I, however, stayed home.

\begin{cousinweb}
\textbf{\mr{मात्र}} (\emph{mātra}) is a tatsama, taken from Sanskrit \emph{mātra}, \textquotedblleft{}measure\textquotedblright{}, which at the end of a compound meant \textquotedblleft{}nothing but, only\textquotedblright{} (Monier-Williams); Marathi uses it on its own for \textquotedblleft{}only\textquotedblright{} and \textquotedblleft{}however\textquotedblright{}.
\end{cousinweb}

\begin{grammarlens}[title={What you've built}]
One more word for giving reasons and opinions.
\end{grammarlens}

\subsection*{Your turn}
\begin{itemize}
\raggedright
  \item \emph{Say it:} \emph{mātra}
  \item \emph{Say it:} \emph{mātra}, once more
  \item \emph{Say it:} \emph{ghar lahān āhe, mātra khūp svaccha āhe.} — The house is small; it is very clean, though.
  \item \emph{Recall:} say the Marathi for to blame, then the Marathi for to accuse, then say \emph{mātra} again
\end{itemize}

\begin{tcolorbox}[breakable,colback=teal!4,colframe=teal!35!black,title={Before you move on}]
What does \mr{मात्र} mean? (\textquotedblleft{}However; though\textquotedblright{}.) Say it once more.
\end{tcolorbox}
\section[\texorpdfstring{\mr{तसेच} (tasec)}{तसेच (tasec)}]{\mr{तसेच} (tasec) — likewise; as well as}

\label{lesson:MR-C351-tasec}



\begin{quote}
Before the new one: say the Marathi for to accuse, then the Marathi for however.
\end{quote}

\subsection*{You'll want to know: \mr{तसेच}}
\textbf{\mr{तसेच}} — \emph{tasec} — \textquotedblleft{}likewise; as well as\textquotedblright{}.

\begin{quote}
\mr{आम्ही} \mr{पुस्तकं} \mr{वाचली}, \mr{तसेच} \mr{गाणी} \mr{ऐकली}.
\end{quote}

\begin{quote}
\emph{āmhī pustakaṁ vāclī, tasec gāṇī aiklī.}
\end{quote}

\begin{quote}
We read books and also listened to songs.
\end{quote}

Formal; \mr{तसंच} in speech. Built from \mr{तसे}, in that way, and -\mr{च}.

\begin{cousinweb}
\textbf{\mr{तसेच}} (\emph{tasec}) is built from \mr{तसे}, \textquotedblleft{}in that way\textquotedblright{}, and the stressing ending -\mr{च}. \mr{तसा} is inherited: CDIAL lists Marathi \emph{tasā}, \textquotedblleft{}suchlike\textquotedblright{}, under Sanskrit \emph{tādṛśa}, \textquotedblleft{}suchlike\textquotedblright{} (CDIAL 5760).
\end{cousinweb}

\begin{grammarlens}[title={What you've built}]
One more word for giving reasons and opinions.
\end{grammarlens}

\subsection*{Your turn}
\begin{itemize}
\raggedright
  \item \emph{Say it:} \emph{tasec}
  \item \emph{Say it:} \emph{tasec}, once more
  \item \emph{Say it:} \emph{āmhī pustakaṁ vāclī, tasec gāṇī aiklī.} — We read books and also listened to songs.
  \item \emph{Recall:} say the Marathi for to accuse, then the Marathi for however, then say \emph{tasec} again
\end{itemize}

\begin{tcolorbox}[breakable,colback=teal!4,colframe=teal!35!black,title={Before you move on}]
What does \mr{तसेच} mean? (\textquotedblleft{}Likewise; as well as\textquotedblright{}.) Say it once more.
\end{tcolorbox}
\section[\texorpdfstring{\mr{किंबहुना} (kimbahunā)}{किंबहुना (kimbahunā)}]{\mr{किंबहुना} (kimbahunā) — in fact; rather}

\label{lesson:MR-C351-kimbahuna}



\begin{quote}
Before the new one: say the Marathi for however, then the Marathi for likewise.
\end{quote}

\subsection*{You'll want to know: \mr{किंबहुना}}
\textbf{\mr{किंबहुना}} — \emph{kimbahunā} — \textquotedblleft{}in fact; rather\textquotedblright{}.

\begin{quote}
\mr{तो} \mr{हुशार} \mr{आहे}, \mr{किंबहुना} \mr{वर्गात} \mr{सर्वात} \mr{हुशार} \mr{आहे}.
\end{quote}

\begin{quote}
\emph{to huśār āhe, kimbahunā vargāt sarvāt huśār āhe.}
\end{quote}

\begin{quote}
He is clever; in fact, the cleverest in the class.
\end{quote}

Formal: it makes what was said stronger.

\begin{cousinweb}
\textbf{\mr{किंबहुना}} (\emph{kimbahunā}) is a tatsama phrase, taken whole from Sanskrit \emph{kiṁ}, \textquotedblleft{}what\textquotedblright{}, and \emph{bahunā}, \textquotedblleft{}with more\textquotedblright{} (Monier-Williams): what need is there to say more.
\end{cousinweb}

\begin{grammarlens}[title={What you've built}]
One more word for giving reasons and opinions.
\end{grammarlens}

\subsection*{Your turn}
\begin{itemize}
\raggedright
  \item \emph{Say it:} \emph{kimbahunā}
  \item \emph{Say it:} \emph{kimbahunā}, once more
  \item \emph{Say it:} \emph{to huśār āhe, kimbahunā vargāt sarvāt huśār āhe.} — He is clever; in fact, the cleverest in the class.
  \item \emph{Recall:} say the Marathi for however, then the Marathi for likewise, then say \emph{kimbahunā} again
\end{itemize}

\begin{tcolorbox}[breakable,colback=teal!4,colframe=teal!35!black,title={Before you move on}]
What does \mr{किंबहुना} mean? (\textquotedblleft{}In fact; rather\textquotedblright{}.) Say it once more.
\end{tcolorbox}
\section[\texorpdfstring{\mr{तथापि} (tathāpi)}{तथापि (tathāpi)}]{\mr{तथापि} (tathāpi) — nevertheless}

\label{lesson:MR-C351-tathapi}



\begin{quote}
Before the new one: say the Marathi for likewise, then the Marathi for in fact.
\end{quote}

\subsection*{You'll want to know: \mr{तथापि}}
\textbf{\mr{तथापि}} — \emph{tathāpi} — \textquotedblleft{}nevertheless\textquotedblright{}.

\begin{quote}
\mr{पाऊस} \mr{पडत} \mr{होता}, \mr{तथापि} \mr{कार्यक्रम} \mr{वेळेवर} \mr{सुरू} \mr{झाला}.
\end{quote}

\begin{quote}
\emph{pāūs paḍat hotā, tathāpi kāryakram veḷevar surū jhālā.}
\end{quote}

\begin{quote}
It was raining; nevertheless, the programme began on time.
\end{quote}

Formal and written; in speech \mr{पण} or \mr{तरीही}.

\begin{cousinweb}
\textbf{\mr{तथापि}} (\emph{tathāpi}) is a tatsama, taken from Sanskrit \emph{tathāpi}, \textquotedblleft{}even so, nevertheless\textquotedblright{} (CDIAL 5647), from \emph{tathā}, \textquotedblleft{}so\textquotedblright{}, and \emph{api}, \textquotedblleft{}even\textquotedblright{}.
\end{cousinweb}

\begin{grammarlens}[title={What you've built}]
One more word for giving reasons and opinions.
\end{grammarlens}

\subsection*{Your turn}
\begin{itemize}
\raggedright
  \item \emph{Say it:} \emph{tathāpi}
  \item \emph{Say it:} \emph{tathāpi}, once more
  \item \emph{Say it:} \emph{pāūs paḍat hotā, tathāpi kāryakram veḷevar surū jhālā.} — It was raining; nevertheless, the programme began on time.
  \item \emph{Recall:} say the Marathi for likewise, then the Marathi for in fact, then say \emph{tathāpi} again
\end{itemize}

\begin{tcolorbox}[breakable,colback=teal!4,colframe=teal!35!black,title={Before you move on}]
What does \mr{तथापि} mean? (\textquotedblleft{}Nevertheless\textquotedblright{}.) Say it once more.
\end{tcolorbox}
\section[\texorpdfstring{\mr{उदाहरणार्थ} (udāharaṇārtha)}{उदाहरणार्थ (udāharaṇārtha)}]{\mr{उदाहरणार्थ} (udāharaṇārtha) — for example}

\label{lesson:MR-C351-udaharanartha}



\begin{quote}
Before the new one: say the Marathi for in fact, then the Marathi for nevertheless.
\end{quote}

\subsection*{You'll want to know: \mr{उदाहरणार्थ}}
\textbf{\mr{उदाहरणार्थ}} — \emph{udāharaṇārtha} — \textquotedblleft{}for example\textquotedblright{}.

\begin{quote}
\mr{मला} \mr{फळं} \mr{आवडतात}, \mr{उदाहरणार्थ} \mr{आंबा} \mr{आणि} \mr{केळं}.
\end{quote}

\begin{quote}
\emph{malā phaḷaṁ āvaḍtāt, udāharaṇārtha āmbā āṇi keḷaṁ.}
\end{quote}

\begin{quote}
I like fruit, for example mangoes and bananas.
\end{quote}

\mr{उदाहरण} + \mr{अर्थ}: for the sake of an example.

\begin{cousinweb}
\textbf{\mr{उदाहरणार्थ}} (\emph{udāharaṇārtha}) is a learned compound built from Sanskrit \emph{udāharaṇa}, \textquotedblleft{}example\textquotedblright{}, and \emph{artha}, \textquotedblleft{}aim, sake\textquotedblright{} (Monier-Williams).
\end{cousinweb}

\begin{grammarlens}[title={What you've built}]
One more word for giving reasons and opinions.
\end{grammarlens}

\subsection*{Your turn}
\begin{itemize}
\raggedright
  \item \emph{Say it:} \emph{udāharaṇārtha}
  \item \emph{Say it:} \emph{udāharaṇārtha}, once more
  \item \emph{Say it:} \emph{malā phaḷaṁ āvaḍtāt, udāharaṇārtha āmbā āṇi keḷaṁ.} — I like fruit, for example mangoes and bananas.
  \item \emph{Recall:} say the Marathi for in fact, then the Marathi for nevertheless, then say \emph{udāharaṇārtha} again
\end{itemize}

\begin{tcolorbox}[breakable,colback=teal!4,colframe=teal!35!black,title={Before you move on}]
What does \mr{उदाहरणार्थ} mean? (\textquotedblleft{}For example\textquotedblright{}.) Say it once more.
\end{tcolorbox}
